\section{Qué cambia si hay ciclos en el grafo local}
Para cualquier colección de grupos, el argumento de conexidad de la parte 3 vale sin
que el problema local de Bethe sea exacto. En cada coeficiente con activaciones
separadas se usa la independencia del exacto y la separación en bloques de los respectivos
predictores. Por tanto, después de alternar, el orden genérico es al menos $s^{c-1}$.
El caso de orden cero es la suma de valores locales y concuerda con esa cota.

Supongamos ahora que en todo el grafo hay un ciclo independiente. Sólo una componente lo
contiene. En el coeficiente de orden $w=c-1$ que alcance los $c$ grupos, necesariamente
$L_I=0$ y todas las interacciones tienen déficit cero. Sólo quedan productos de formas
lineales de acción, cada una sobre un grupo distinto. Se enraíza el árbol en la componente
con el ciclo; los demás grupos son forestales. La eliminación de ramas de la parte 4
reduce el coeficiente a derivadas en campos locales de la diferencia de logaritmos de la
raíz, multiplicadas por funciones holomorfas acotadas de las hojas. Incluye $c=1,w=0$
sin necesidad de eliminación. Las derivadas en campos son combinaciones de
$\zeta_i\partial_{\zeta_i}$, pues una actividad se multiplica al añadirle un campo.

En la raíz, en cada activación, el ciclo se rompe o se conserva con los mismos sitios.
Al quitar factores o incidencias jamás se crea otro ciclo. En el primer caso la diferencia
local es cero. En el segundo, el cociente de partición exacta respecto a Bethe es
$1+\theta$, y el prefactor del predictor es $1/(1-\theta)$, con
\[
 \theta=(-1)^b\prod_{i\text{ en el ciclo}} m_i/(1-m_i).
\]
Aquí $b$ es el número de sitios del ciclo y $m_i$ son las creencias de Bethe en sus campos.
Son las identidades local binaria y gaussiana de la parte 4 del estudio interlíneas.
El error $\log(1-\theta^2)$ tiene dos factores $\zeta_i$ por cada sitio del ciclo:
$m_i$ es divisible por $\zeta_i$ y el denominador es una unidad holomorfa.
Las derivadas en campos preservan esas potencias. En particular el término de orden
$s^{c-1}$ gana al menos una potencia física adicional incluso después de las activaciones.

\paragraph{Uso de las caras con un tamaño entero.}
Para $j$ fijo se usa el polidisco $|\zeta_i|\le\rho_v$, con $\rho_v>0$ independiente de $j$.
El valor exacto es el logaritmo de una suma finita próxima a uno. La silla y el determinante
regularizado son holomorfos y acotados tras sustraer las constantes y campos comunes.
Este control conjunto en las actividades es el de las partes 2--3 del estudio interlíneas;
para los órdenes superiores se usan los Taylor uniformes descritos en las partes 2 y 3.
En una cara $\zeta_i=0$ el caso que activa $i$ y el que lo desactiva coinciden,
también para el determinante. Por ello $H=[S]e$ satisface
$H=\mathcal U H$, con $\mathcal U=\prod_{i\in S}(I-\mathrm{ev}_{\zeta_i=0})$.
Se aplica $\mathcal U$ por separado a los coeficientes y al resto de cualquier
descomposición uniforme de $H$ para ese entero $j$. Su norma es a lo sumo $2^v$;
elimina los monomios a los que falta algún sitio y conserva la potencia doble de los
sitios del ciclo. Se gana así el producto de $v$ actividades mediante Cauchy en un
polidisco menor. Se evalúa al final en $|\zeta_i|\le C_v s$. No se exige que los restos
de la expansión física sean diferenciables en el tamaño del tablero.
Para el peso individual $|[S]e|\le C_v j^{-p}$ resulta
\[
 p=\begin{cases}
 v+c+1 & \text{si todas las componentes son bosques},\\
 v+c   & \text{si hay exactamente un ciclo independiente en total},\\
 v+c-1 & \text{en los demás casos}.
 \end{cases}
\]
En el segundo caso el resto de orden $s^c$ ya da $s^{v+c}$; para el coeficiente
anterior se usan los factores del ciclo además de la proyección.
